\documentclass[11pt,a4paper]{article}

\usepackage[utf8]{inputenc}
\usepackage[T1]{fontenc}
\usepackage[french]{babel}
\usepackage{geometry}
\usepackage{booktabs}
\usepackage{hyperref}
\usepackage{amsmath}
\usepackage{microtype}
\usepackage{listings}
\usepackage{graphicx}
\usepackage{xcolor}
\usepackage{tikz}
\usepackage{tcolorbox}
\usepackage{tabularx}
\usepackage{array}
\usepackage{enumitem}
\usepackage{mdframed}
\usepackage{fancyhdr}
\usepackage{titlesec}
\usepackage{multicol}

\usetikzlibrary{arrows.meta, shapes.geometric, positioning, fit, backgrounds}

\geometry{margin=2.5cm, top=3cm, bottom=3cm}

% ---------- couleurs projet ----------
\definecolor{brandblue}{RGB}{30, 90, 160}
\definecolor{brandgray}{RGB}{240, 242, 246}
\definecolor{warnred}{RGB}{200, 50, 50}
\definecolor{okgreen}{RGB}{34, 139, 80}
\definecolor{codebg}{RGB}{248, 248, 252}
\definecolor{accentorange}{RGB}{220, 120, 30}

% ---------- en-têtes ----------
\pagestyle{fancy}
\fancyhf{}
\rhead{\small\textcolor{brandblue}{IA Personnalisation Produit — MVP}}
\lhead{\small\textcolor{gray}{Revue de mi-parcours}}
\rfoot{\small\thepage}
\renewcommand{\headrulewidth}{0.4pt}

% ---------- titres de sections ----------
\titleformat{\section}{\large\bfseries\color{brandblue}}{}{0em}{\thesection\quad}[\titlerule]
\titleformat{\subsection}{\normalsize\bfseries\color{brandblue!80!black}}{}{0em}{\thesubsection\quad}

% ---------- boîtes ----------
\tcbuselibrary{skins, breakable}

\newtcolorbox{infobox}[1]{
  colback=brandgray, colframe=brandblue,
  fonttitle=\bfseries\small, title=#1,
  left=4pt, right=4pt, top=4pt, bottom=4pt,
  breakable
}

\newtcolorbox{warnbox}[1]{
  colback=warnred!6, colframe=warnred,
  fonttitle=\bfseries\small\color{warnred}, title=#1,
  left=4pt, right=4pt, top=4pt, bottom=4pt,
  breakable
}

\newtcolorbox{okbox}[1]{
  colback=okgreen!6, colframe=okgreen,
  fonttitle=\bfseries\small\color{okgreen}, title=#1,
  left=4pt, right=4pt, top=4pt, bottom=4pt,
  breakable
}

% ---------- style listings ----------
\lstset{
  backgroundcolor=\color{codebg},
  basicstyle=\ttfamily\small,
  breaklines=true,
  frame=single,
  framerule=0.4pt,
  rulecolor=\color{brandblue!40},
  keywordstyle=\color{brandblue}\bfseries,
  stringstyle=\color{okgreen},
  commentstyle=\color{gray}\itshape,
  numbers=left,
  numberstyle=\tiny\color{gray},
  numbersep=8pt,
  xleftmargin=12pt,
  tabsize=2,
  showstringspaces=false
}

% ---------- commandes utiles ----------
\newcommand{\badge}[2]{%
  \colorbox{#1}{\textcolor{white}{\footnotesize\textbf{~#2~}}}%
}

% ======================================================
\title{%
  \vspace{-1cm}
  \textbf{\color{brandblue}Revue de mi-parcours}\\[4pt]
  \large IA de personnalisation produit — MVP\\[6pt]
  \normalsize\textcolor{gray}{Génération de contenus sous contraintes réglementaires}
}
\author{}
\date{\textcolor{gray}{\small Mars 2026}}
% ======================================================

\begin{document}

\maketitle
\thispagestyle{fancy}

\tableofcontents
\vspace{1em}

% ======================================================
\section{Résumé exécutif}

\begin{infobox}{En un coup d'œil}
Ce projet automatise la génération de contenus produits pour la santé animale dans un cadre réglementé.
Le MVP repose sur un pipeline \textbf{contrôlé, structuré et traçable} capable de produire des contenus
conformes, cohérents avec le produit et adaptés au canal de distribution.
\end{infobox}

\vspace{0.5em}

\begin{center}
\begin{tabular}{lll}
\toprule
\textbf{Dimension} & \textbf{État} & \textbf{Note} \\
\midrule
Fonctionnalité de base   & \badge{okgreen}{OK}      & Pipeline bout-en-bout opérationnel \\
Conformité réglementaire & \badge{okgreen}{OK}      & Pare-feu sémantique + règles déterministes \\
Traçabilité              & \badge{okgreen}{OK}      & Logs JSON complets par exécution \\
Passage à l'échelle      & \badge{accentorange}{En cours} & Batch et multi-retry à implémenter \\
Validation humaine       & \badge{accentorange}{Partiel}  & Interface à industrialiser \\
\bottomrule
\end{tabular}
\end{center}

% ======================================================
\section{Problématique}

\subsection{Contraintes métier}

Les produits de santé animale sont soumis à des restrictions strictes sur les allégations,
définies par les réglementations européennes (directive 2001/82/CE, règlement 2019/6) et les
recommandations DGCCRF pour les compléments alimentaires vétérinaires.

\begin{center}
\renewcommand{\arraystretch}{1.3}
\begin{tabularx}{\linewidth}{>{\bfseries}lX}
\toprule
Statut & Exemples d'allégations \\
\midrule
\badge{warnred}{Interdit} &
  \textit{guérit}, \textit{traite}, \textit{élimine}, \textit{garanti},
  \textit{médicament}, \textit{thérapeutique} \\
\badge{okgreen}{Autorisé} &
  \textit{aide à}, \textit{soutient}, \textit{contribue à}, \textit{favorise},
  \textit{participe au bien-être de} \\
\badge{accentorange}{Conditionnel} &
  \textit{réduit} (si données cliniques disponibles et citées),
  allégations nutritionnelles validées \\
\bottomrule
\end{tabularx}
\end{center}

\subsection{Tension structurelle}

Un LLM non contraint optimise naturellement la \textit{persuasion}, au détriment de la
\textit{précision réglementaire}. Cette tension est fondamentale :

\begin{multicols}{2}
\begin{warnbox}{Texte marketing non contraint}
\small
\textit{« Ce complément \textbf{guérit} efficacement les troubles digestifs de votre chat et
\textbf{élimine garantie} les inconforts intestinaux. »}
\end{warnbox}

\columnbreak

\begin{okbox}{Texte conforme généré}
\small
\textit{« Ce complément \textbf{aide à soutenir} l'équilibre de la flore intestinale et
\textbf{contribue au confort digestif} de votre chat. »}
\end{okbox}
\end{multicols}

\subsection{Objectif}

Générer automatiquement des contenus :

\begin{itemize}[leftmargin=1.5em]
  \item \textbf{conformes réglementairement} — aucune allégation interdite, longueurs
        respectées, mentions légales présentes ;
  \item \textbf{fidèles au produit} — basés sur la fiche technique officielle du produit ;
  \item \textbf{adaptés au canal} — ton et format différenciés (e-commerce, vétérinaire,
        distribution spécialisée).
\end{itemize}

% ======================================================
\section{Approche}

Le principe central est de \textbf{contrôler l'espace de génération en amont} plutôt que de
corriger a posteriori. Chaque étape produit des artefacts vérifiables et versionnés.

\subsection{Pipeline général}

\begin{center}
\begin{tikzpicture}[
  node distance=1.1cm and 0.4cm,
  box/.style={
    rectangle, rounded corners=4pt,
    minimum width=2.2cm, minimum height=0.8cm,
    text centered, font=\small\bfseries,
    draw, thick
  },
  step/.style={box, fill=brandblue!15, draw=brandblue},
  judge/.style={box, fill=accentorange!20, draw=accentorange},
  human/.style={box, fill=okgreen!15, draw=okgreen},
  arr/.style={-{Stealth[length=5pt]}, thick, brandblue},
  retarr/.style={-{Stealth[length=5pt]}, thick, warnred, dashed}
]

\node[step]  (ext)  {extract};
\node[step, right=of ext]  (ana)  {analyze};
\node[step, right=of ana]  (gen)  {generate};
\node[judge, right=of gen] (jud)  {judge};
\node[step, right=of jud]  (val)  {validate};
\node[human, right=of val] (hum)  {human-in\\-the-loop};

\draw[arr] (ext) -- (ana);
\draw[arr] (ana) -- (gen);
\draw[arr] (gen) -- (jud);
\draw[arr] (jud) -- (val);
\draw[arr] (val) -- (hum);

% boucle retry
\draw[retarr] (jud.south) -- ++(0,-0.6) -| (gen.south)
  node[midway, below, font=\footnotesize\color{warnred}] {retry si non-conforme};

% annotations
\node[font=\tiny, text=gray, above=0.1cm of ext] {CSV produit};
\node[font=\tiny, text=gray, above=0.1cm of ana] {brief + contraintes};
\node[font=\tiny, text=gray, above=0.1cm of gen] {contenus bruts};
\node[font=\tiny, text=gray, above=0.1cm of jud] {scores};
\node[font=\tiny, text=gray, above=0.1cm of val] {rapport};
\node[font=\tiny, text=gray, above=0.1cm of hum] {décision finale};

\end{tikzpicture}
\end{center}

\subsection{Principe de contrôle en amont}

\[
\underbrace{\text{extract} \to \text{analyze}}_{\text{grounding}} \to
\underbrace{\text{generate}}_{\text{génération contrainte}} \to
\underbrace{\text{judge} \to \text{validate}}_{\text{double vérification}} \to
\underbrace{\text{human}}_{\text{supervision}}
\]

Le \textit{grounding} est l'étape clé : il transforme une fiche technique brute en un ensemble
de contraintes explicites qui guident (et limitent) la génération.

% ======================================================
\section{Description du système}

\subsection{Extraction (\texttt{extract.py})}

Lecture et normalisation du fichier source CSV. Chaque ligne correspond à un produit ;
les colonnes sont standardisées en un objet produit unifié.

\begin{lstlisting}[language=Python, caption={Exemple de sortie extract}]
{
  "product_id": "CAT-CALM-001",
  "name": "ZenCat Relax",
  "species": "chat",
  "category": "complement_comportement",
  "active_ingredients": ["Valériane", "L-Tryptophane", "Magnésium"],
  "form": "comprimé",
  "dosage": "1 comprimé / 10 kg / jour"
}
\end{lstlisting}

\subsection{Analyse et grounding (\texttt{analyze.py})}

Construction d'un \textbf{brief produit structuré} et d'un ensemble de contraintes de conformité
à partir des données extraites. Cette étape est le cœur du contrôle réglementaire.

\begin{lstlisting}[language=Python, caption={Brief de conformité produit}]
{
  "product_id": "CAT-CALM-001",
  "allowed_claims": [
    "aide à réduire le stress situationnel",
    "contribue au calme naturel",
    "soutient l'équilibre émotionnel"
  ],
  "forbidden_claims": [
    "guérit l'anxiété",
    "traite les troubles comportementaux",
    "sédatif garanti"
  ],
  "mandatory_mentions": [
    "complément alimentaire",
    "ne se substitue pas à un traitement vétérinaire"
  ],
  "channel_tone": {
    "ecommerce": "chaleureux, bénéfices consommateur",
    "veterinaire": "factuel, ingrédients actifs mis en avant"
  }
}
\end{lstlisting}

\begin{infobox}{Pare-feu sémantique}
Le brief de conformité agit comme un \textbf{contexte de contraintes injecté dans chaque prompt}.
Le modèle ne peut pas ignorer les contraintes sans produire un résultat explicitement incohérent,
ce qui facilite leur détection au stade de l'évaluation.
\end{infobox}

\subsection{Génération (\texttt{generate.py})}

Production des contenus pour chaque canal et format :

\begin{center}
\renewcommand{\arraystretch}{1.3}
\begin{tabular}{lll}
\toprule
\textbf{Format} & \textbf{Longueur cible} & \textbf{Usage} \\
\midrule
Description courte  & 150–250 car. & Listing, résumé \\
Description moyenne & 400–600 car. & Fiche produit principale \\
Description longue  & 800–1200 car. & Page détaillée, blog \\
SEO title           & 50–60 car.   & Balise \texttt{<title>} \\
Meta description    & 150–160 car. & Balise \texttt{<meta>} \\
Arguments commerciaux & 3–5 puces  & Vendeur, distributeur \\
\bottomrule
\end{tabular}
\end{center}

\subsection{Évaluation (\texttt{judge.py})}

Un LLM juge évalue chaque contenu généré selon trois axes, avec un score normalisé et
une justification structurée :

\begin{center}
\renewcommand{\arraystretch}{1.3}
\begin{tabular}{llp{7cm}}
\toprule
\textbf{Axe} & \textbf{Score} & \textbf{Critères} \\
\midrule
Conformité réglementaire & 0–10 &
  Absence d'allégations interdites, présence des mentions obligatoires \\
Qualité SEO & 0–10 &
  Densité de mots-clés, longueur, structure, lisibilité \\
Cohérence produit & 0–10 &
  Fidélité aux ingrédients, à l'espèce cible, aux bénéfices documentés \\
\midrule
\textbf{Score global} & \textbf{0–10} & Moyenne pondérée (conformité $\times$ 2) \\
\bottomrule
\end{tabular}
\end{center}

\noindent Un seuil de validation est fixé : \textbf{conformité $\geq$ 8 et score global $\geq$ 7}.
Tout contenu en dessous déclenche une boucle de correction.

\subsection{Correction (\texttt{retry.py})}

La correction est \textbf{ciblée} : seuls les axes dont le score est insuffisant sont re-générés.
Le prompt de correction injecte explicitement les éléments défaillants identifiés par le juge.

\begin{lstlisting}[language=Python, caption={Payload de correction ciblée}]
{
  "content_id": "desc_courte",
  "failed_axes": ["conformite"],
  "feedback": "Allégation interdite détectée : 'réduit les crises d'anxiété'.
    Remplacer par une formulation du type 'aide à apaiser'.",
  "original_text": "..."
}
\end{lstlisting}

\subsection{Validation déterministe}

En complément de l'évaluation LLM, des règles déterministes garantissent un dernier
niveau de vérification non probabiliste :

\begin{itemize}[leftmargin=1.5em]
  \item \textbf{Regex} — liste noire de termes interdits (mise à jour manuelle) ;
  \item \textbf{Longueur} — contrôle des bornes min/max par format ;
  \item \textbf{SEO} — vérification de la densité de mots-clés cibles ;
  \item \textbf{Mentions légales} — présence obligatoire des disclaimers.
\end{itemize}

\begin{warnbox}{Complémentarité des approches}
L'évaluation LLM capte les violations sémantiques subtiles ; la validation déterministe
garantit les invariants non ambigus. Les deux couches sont nécessaires : ni l'une ni l'autre
n'est suffisante seule.
\end{warnbox}

\subsection{Human-in-the-loop}

Dans un contexte réglementé, la validation humaine est une exigence, non une option.
Elle s'insère après la validation déterministe et avant toute publication.

\begin{center}
\begin{tikzpicture}[
  node distance=0.8cm,
  box/.style={rectangle, rounded corners=3pt, draw, thick,
              minimum width=3cm, minimum height=0.7cm,
              font=\small, text centered},
  dec/.style={diamond, draw, thick, fill=brandgray,
              font=\small, text centered, aspect=3,
              minimum width=3.5cm},
  arr/.style={-{Stealth[length=4pt]}, thick}
]
\node[box, fill=brandblue!15, draw=brandblue] (gen) {Contenu généré + scores};
\node[dec, below=of gen] (rev) {Révision humaine};
\node[box, fill=okgreen!15, draw=okgreen, below left=0.8cm and 1cm of rev]  (ok)  {Validé → export};
\node[box, fill=accentorange!15, draw=accentorange, below=0.8cm of rev]     (mod) {Modifié → export};
\node[box, fill=warnred!15, draw=warnred, below right=0.8cm and 1cm of rev] (rej) {Rejeté → feedback};

\draw[arr] (gen) -- (rev);
\draw[arr] (rev) -- node[left,font=\tiny]{valide} (ok);
\draw[arr] (rev) -- node[right,font=\tiny]{modifie} (mod);
\draw[arr] (rev) -- node[right,font=\tiny]{rejette} (rej);

\draw[arr, dashed, warnred] (rej.east) -- ++(1.2,0) |- (gen.east)
  node[midway, right, font=\tiny\color{warnred}] {amélioration continue};
\end{tikzpicture}
\end{center}

Les retours humains (modifications, rejets, commentaires) constituent un corpus d'apprentissage
pour l'amélioration continue des prompts et des critères de jugement.

% ======================================================
\section{Implémentation}

\subsection{Structure des modules}

\begin{center}
\renewcommand{\arraystretch}{1.3}
\begin{tabular}{lll}
\toprule
\textbf{Module} & \textbf{Responsabilité} & \textbf{Entrée / Sortie} \\
\midrule
\texttt{extract.py}  & Ingestion et normalisation & CSV $\to$ \texttt{Product} \\
\texttt{analyze.py}  & Grounding réglementaire    & \texttt{Product} $\to$ \texttt{Brief} \\
\texttt{generate.py} & Génération LLM             & \texttt{Brief} $\to$ \texttt{ContentSet} \\
\texttt{judge.py}    & Évaluation multi-axes      & \texttt{ContentSet} $\to$ \texttt{Scores} \\
\texttt{retry.py}    & Correction ciblée          & \texttt{ContentSet} + \texttt{Scores} $\to$ \texttt{ContentSet} \\
\bottomrule
\end{tabular}
\end{center}

\subsection{Traçabilité}

Chaque exécution produit un fichier de log JSON structuré, horodaté et versionné :

\begin{lstlisting}[language=Python, caption={Structure d'un log d'exécution}]
{
  "run_id": "2026-03-27T09:14:22Z_CAT-CALM-001",
  "product_id": "CAT-CALM-001",
  "pipeline_version": "0.4.1",
  "steps": {
    "extract":  { "status": "ok", "duration_ms": 12 },
    "analyze":  { "status": "ok", "duration_ms": 1840, "model": "claude-sonnet-4-6" },
    "generate": { "status": "ok", "duration_ms": 4320, "tokens_used": 1247 },
    "judge":    { "status": "ok", "scores": { "conformite": 9.1,
                                               "seo": 7.8, "coherence": 8.5 } },
    "retry":    { "triggered": false },
    "validate": { "status": "ok", "regex_hits": 0 }
  },
  "prompts": { "analyze": "...", "generate": "...", "judge": "..." }
}
\end{lstlisting}

La conservation des prompts dans les logs est essentielle pour la reproductibilité et
l'amélioration itérative des templates.

\subsection{Stack technique}

\begin{multicols}{2}
\begin{itemize}[leftmargin=1em, itemsep=2pt]
  \item Python 3.11
  \item Claude API (Anthropic SDK)
  \item Pydantic — validation des schémas
  \item \texttt{re} — validation déterministe
  \item JSON Lines — logs structurés
  \item pytest — tests unitaires des règles
\end{itemize}
\end{multicols}

% ======================================================
\section{Limites actuelles}

\begin{warnbox}{Limites identifiées}
\begin{enumerate}[leftmargin=1.5em, itemsep=4pt]
  \item \textbf{Pas de batch} — le pipeline traite un produit à la fois ;
        les CSV de plusieurs centaines de références ne sont pas gérables efficacement.
  \item \textbf{Un seul cycle de correction} — si le premier retry ne suffit pas,
        le contenu est marqué comme non résolu sans nouvelle tentative.
  \item \textbf{Dépendance stricte au CSV} — tout changement de format de la source
        peut casser le pipeline sans message d'erreur explicite.
  \item \textbf{Cohérence globale non garantie} — deux descriptions générées
        indépendamment pour le même produit peuvent présenter des contradictions
        de ton ou de formulation.
  \item \textbf{Coût LLM non optimisé} — chaque étape (analyze, generate, judge)
        effectue des appels indépendants ; le caching de contexte n'est pas exploité.
\end{enumerate}
\end{warnbox}

% ======================================================
\section{Prochaines étapes}

\begin{center}
\renewcommand{\arraystretch}{1.3}
\begin{tabularx}{\linewidth}{clXl}
\toprule
\textbf{Priorité} & \textbf{Fonctionnalité} & \textbf{Description} & \textbf{Horizon} \\
\midrule
P1 & Batch processing &
  Traitement parallèle de N produits avec gestion des erreurs par produit &
  T+1 mois \\
P1 & Multi-retry avec convergence &
  Boucle de correction avec seuil de sortie configurable (max 3 cycles) &
  T+1 mois \\
P2 & Validation humaine industrialisée &
  Interface légère (web ou CLI) pour review, modification et feedback &
  T+2 mois \\
P2 & Cohérence inter-descriptions &
  Prompt de cohérence qui prend en entrée l'ensemble des variantes &
  T+2 mois \\
P3 & Optimisation des coûts &
  Prompt caching, modèles légers pour les étapes déterministes &
  T+3 mois \\
P3 & Monitoring en production &
  Dashboard des scores, alertes sur dérives de conformité &
  T+3 mois \\
\bottomrule
\end{tabularx}
\end{center}

% ======================================================
\section{Conclusion}

\begin{okbox}{Bilan}
Le MVP démontre la \textbf{faisabilité d'une génération contrôlée sous contraintes
réglementaires strictes}. Le système est fonctionnel, traçable et modulaire.

Les prochaines étapes prioritaires sont le \textbf{batch} et la \textbf{boucle de convergence},
qui sont les conditions nécessaires à un déploiement industriel.
\end{okbox}

\vspace{0.5em}

La principale valeur architecturale du MVP est la séparation claire entre :
\begin{itemize}[leftmargin=1.5em]
  \item le \textit{grounding} (ce que le modèle \textbf{peut} dire),
  \item la \textit{génération} (ce que le modèle \textbf{dit}),
  \item et la \textit{vérification} (ce que le modèle \textbf{a dit est-il conforme}).
\end{itemize}

Cette séparation est ce qui rend le système \textbf{auditable} — une exigence non négociable
dans un secteur réglementé.

% ======================================================
\appendix
\section{Exemple de sortie complète}

\begin{lstlisting}[caption={ContentSet — ZenCat Relax, canal e-commerce}]
{
  "product_id": "CAT-CALM-001",
  "channel": "ecommerce",
  "desc_courte": "ZenCat Relax aide votre chat à traverser sereinement
    les situations stressantes du quotidien grâce à une formule naturelle
    à base de Valériane et L-Tryptophane. Complément alimentaire.",
  "desc_moyenne": "ZenCat Relax est un complément alimentaire formulé
    pour soutenir l'équilibre émotionnel des chats adultes. Sa composition
    associe la Valériane, traditionnellement reconnue pour contribuer
    au calme naturel, le L-Tryptophane, précurseur de la sérotonine,
    et le Magnésium. Idéal avant un déménagement, une visite vétérinaire
    ou l'arrivée d'un nouvel animal. Ne se substitue pas à un traitement
    vétérinaire.",
  "seo_title": "ZenCat Relax — complément calme et bien-être chat",
  "meta_desc": "Aidez votre chat à rester serein face au stress avec
    ZenCat Relax, complément naturel à la Valériane. Sans ordonnance.",
  "arguments": [
    "Formule naturelle : Valériane, L-Tryptophane, Magnésium",
    "Aide à apaiser les comportements liés au stress situationnel",
    "Palatabilité optimisée — acceptation facilitée",
    "Fabriqué en France, contrôle qualité renforcé"
  ],
  "scores": { "conformite": 9.4, "seo": 8.1, "coherence": 9.0 }
}
\end{lstlisting}

\end{document}
